\section{Discussion}

The results reveal no single dominant configuration, and instead each approach demonstrates distinct trade-offs between plan generation, refusal accuracy, and token cost. The iterative approach achieves the highest rate of successful plan generation on solvable problems and the lowest token cost, as its lower bound is a single LLM query. However, this comes at a cost: the looping structure prevents the system from refusing unsolvable problems under thinking-enabled models, which continue to search for a substitute even when none exists. Applications where false positives are dangerous should favor consensus or judge-based approaches, which preserve refusal accuracy with little sacrifice in plan generation performance.

Increasing the consensus sample count from $n{=}3$ to $n{=}5$ improves performance in some cases --- most notably on bulk problems where additional samples increase the chance of a correct majority vote --- but at roughly double the token cost. Whether this trade-off is worthwhile depends on the target domain; for problems where a single substitution is sufficient, $n{=}3$ is generally adequate.

Model selection has a meaningful impact beyond raw accuracy. Thinking-enabled models (\textsc{DS}, \textsc{G4}) are more willing to make creative type assignments, which benefits solvable problems but undermines refusal on unsolvable ones. \textsc{G4NT} presents a different profile: more conservative and literal, with comparable plan generation accuracy but substantially better refusal behaviour. This suggests that model choice should be informed by the application's tolerance for invalid substitutions. An interesting direction for future work would be a mixture-of-models consensus, combining the creative reasoning of thinking-enabled models with the conservatism of \textsc{G4NT} to improve both plan generation and refusal simultaneously.
